L’animazione digitale rappresenta uno dei campi in cui la tecnologia mostra con maggiore evidenza la propria capacità di simulare il movimento e la realtà. Dietro ogni modello 3D animato si trovano concetti e strumenti propri dell’informatica, come algoritmi, calcoli geometrici e trasformazioni matematiche che consentono di rappresentare movimenti credibili e coerenti nello spazio tridimensionale.

L’animazione di un modello digitale richiede una serie di passaggi fondamentali che consentono di trasformare un oggetto tridimensionale statico in un elemento capace di muoversi in modo credibile. Si parte dalla definizione di una struttura interna che ne guida i movimenti, per poi stabilire la relazione tra questa struttura e la superficie del modello. Infine, vengono applicate le sequenze di movimento che danno vita all’oggetto all’interno della scena virtuale. Queste fasi, pur diverse tra loro, cooperano per rendere il comportamento del modello coerente e realistico, permettendo di ottenere animazioni fluide e naturali.

Negli ultimi anni, l’introduzione di tecniche di automazione e di algoritmi avanzati ha reso possibile generare e trasferire animazioni in modo sempre più efficiente. In questo contesto, il presente lavoro si concentra sull’analisi e l’implementazione di metodi per la gestione automatica del rigging e dello skinning, con l’obiettivo di estendere tali processi anche a modelli non convenzionali e oggetti di natura non umanoide.

L’obiettivo generale è quello di esplorare i principi teorici e le soluzioni pratiche che permettono di applicare in modo coerente un movimento realistico a qualsiasi modello tridimensionale, evidenziando le potenzialità e i limiti delle attuali tecniche di animazione automatizzata. La stesura di questa tesi, scritta in LateX con il metodo top-down, é articolata come segue: 
\begin {enumerate}
    \item Il primo capitolo introduce i concetti fondamentali della rappresentazione e dell’animazione di modelli tridimensionali. Vengono illustrati i principi di costruzione delle mesh, la gestione di vertici e facce e le principali trasformazioni geometriche nello spazio 3D. Si approfondiscono inoltre i meccanismi di rigging, skinning e animazione scheletrica, ponendo le basi per comprendere le fasi operative successive.
    \item Nel secondo capitolo vengono descritti gli strumenti e le tecnologie utilizzate per la realizzazione del progetto. Si analizzano le caratteristiche dei software di modellazione e animazione, i formati di scambio e le modalità di integrazione tra ambienti diversi, delineando il contesto tecnico all’interno del quale si è svolto il lavoro sperimentale.
    \item Il terzo capitolo illustra l’approccio seguito per la creazione e l’animazione dei modelli tridimensionali. Vengono descritte le fasi di generazione automatica dello scheletro, l’associazione delle mesh al rig e il processo di trasferimento delle animazioni. La sezione analizza inoltre le scelte metodologiche adottate e le soluzioni sviluppate per ottimizzare il workflow di animazione.
    \item Il quarto capitolo è dedicato alla parte sperimentale. Qui vengono presentate le applicazioni pratiche delle tecniche descritte, con particolare attenzione ai risultati ottenuti su modelli non convenzionali. L’analisi dei test evidenzia le prestazioni del sistema, i principali limiti riscontrati e le potenzialità dell’approccio proposto.
    \item L’ultimo capitolo raccoglie le conclusioni del lavoro, riassumendo gli obiettivi raggiunti e le criticità incontrate. Vengono inoltre proposte possibili evoluzioni del progetto e nuovi ambiti di applicazione, con l’intento di migliorare ulteriormente l’automazione e la qualità dell’animazione digitale.
\end{enumerate}